\documentclass[11pt]{article}
\usepackage[margin=1in]{geometry}
\usepackage{amsmath,amssymb,amsthm}
\usepackage{booktabs}
\usepackage{hyperref}

\newtheorem{theorem}{Theorem}
\newtheorem{lemma}[theorem]{Lemma}
\newtheorem{proposition}[theorem]{Proposition}
\theoremstyle{remark}
\newtheorem{remark}[theorem]{Remark}

\title{Disproof of Conjecture 00000001102:\\
the R-polynomials of $S_3$ do not peak at $2^{d}(1-2^{-\lceil d/2\rceil})$}
\author{jilint777}
\date{2026-10-04}

\begin{document}
\sloppy
\maketitle

\begin{abstract}
Conjecture 00000001102 asserts that for all $x\le w$ in a Coxeter group, the
Kazhdan--Lusztig R-polynomial satisfies
$\max R_{x,w}=2^{d}\,(1-2^{-\lceil d/2\rceil})$ with $d=\ell(w)-\ell(x)$.
In the symmetric group $S_3$ (type $A_2$) the R-polynomials are
$1,\ q-1,\ (q-1)^2,\ q^3-2q^2+2q-1$ for $d=0,1,2,3$, while the conjectured values
are $0,1,2,6$. We show that the claim fails under every reasonable reading of
``$\max R_{x,w}$'': the largest coefficient, the largest absolute value of a
coefficient, the sum of the absolute values of the coefficients, the leading
coefficient, the degree, the number of nonzero terms, the value at any fixed point
$q_0$, and the maximum of $q\mapsto R_{x,w}(q)$. It fails even for strict pairs
$x<w$ only, and also when ``max'' is taken over all pairs with a given $d$.
The recursion printed in the statement, read literally, forces $R\equiv0$, so it
cannot define anything else. Lean~4 (core library only) formalizes $S_3$, its
Bruhat order, the R-polynomials (characterized by their defining recursion,
with existence and uniqueness proved) and all of these refutations.
\end{abstract}

\section{The conjecture}

The statement (English version) reads:
\begin{quote}
\emph{Definition:} The R-polynomials of KL theory are defined by the recursion
$R_{x,w}=qR_{xs,w}+R_{x,ws}$. \emph{Conjecture:} For all $x\le w$,
$\max R_{x,w}=2^{\ell(w)-\ell(x)}\cdot(1-2^{-\lceil(\ell(w)-\ell(x))/2\rceil})$;
that is, the peak decays as a power of two, with decay exponent the ceiling of
half the length difference.
\end{quote}
The Chinese version says the same. Write $d=\ell(w)-\ell(x)$ and
\[
f(d)=2^{d}\bigl(1-2^{-\lceil d/2\rceil}\bigr)=2^d-2^{d-\lceil d/2\rceil}=2^d-2^{\lfloor d/2\rfloor},
\]
so $f(0)=0$, $f(1)=1$, $f(2)=2$, $f(3)=6$, $f(4)=12$.

\paragraph{Which R-polynomials.} ``The R-polynomials of KL theory'' are the
polynomials $R_{x,w}\in\mathbb Z[q]$ of Kazhdan and Lusztig~\cite{kl}, defined
for a Coxeter system $(W,S)$ by
\[
T_{w^{-1}}^{-1}=\varepsilon_w q^{-\ell(w)}\sum_{x\le w}\varepsilon_x R_{x,w}(q)\,T_x
\]
in the Hecke algebra ($\varepsilon_x=(-1)^{\ell(x)}$). They are characterized
by (see \cite[\S7.5]{humphreys}, \cite[Ch.~5]{bb}):
\begin{equation}\label{eq:R}
\begin{aligned}
&R_{x,w}=0 \text{ unless } x\le w,\qquad R_{w,w}=1,\\
&\text{if } s\in S,\ ws<w: \quad R_{x,w}=\begin{cases}R_{xs,ws}& xs<x,\\
(q-1)R_{x,ws}+qR_{xs,ws}& xs>x.\end{cases}
\end{aligned}
\end{equation}
The recursion printed in the statement, $R_{x,w}=qR_{xs,w}+R_{x,ws}$, has no base
case and no condition on $s$; it is a garbled version of~\eqref{eq:R}. Read
literally (for all $x,w$ and one fixed $s$) it forces $R\equiv0$
(Proposition~\ref{prop:literal}), which contradicts $R_{w,w}=1$; then
$\max R_{x,w}=0\ne1=f(1)$ for every covering pair anyway. We therefore use the
standard R-polynomials~\eqref{eq:R}.

\paragraph{Which group.} The statement concerns KL theory in general, with no
particular Coxeter group named. The symmetric group $S_3$ (type $A_2$) is a Coxeter
group, so a counterexample in $S_3$ refutes ``for all $x\le w$''.

\paragraph{Which ``max''.} ``$\max R_{x,w}$'' (the ``peak'') is not defined in the
statement. We consider every reading we could think of
(Section~\ref{sec:readings}) and refute each one, already restricted to strict
pairs $x<w$. The pair $x=w$, where $R_{w,w}=1$ but $f(0)=0$, gives a second
counterexample for every reading $\rho$ with $\rho(1)\ne0$.

\section{\texorpdfstring{$S_3$, its Bruhat order and its R-polynomials}{S3, its Bruhat order and its R-polynomials}}

We write permutations of $\{1,2,3\}$ in one-line notation $w=[w(1),w(2),w(3)]$.
The length $\ell(w)$ is the number of inversions. The reflections are the three
transpositions; right multiplication $w\mapsto wt$ by the transposition $t=(i\,j)$
swaps the entries in positions $i$ and~$j$. The simple reflections are
$s_1=(1\,2)$ and $s_2=(2\,3)$. Bruhat order is the reflexive--transitive closure
of $x\to xt$ ($t$ a reflection, $\ell(xt)>\ell(x)$). Write
$e=[1,2,3]$, $s_1=[2,1,3]$, $s_2=[1,3,2]$, $s_1s_2=[2,3,1]$, $s_2s_1=[3,1,2]$,
$w_0=[3,2,1]$. The Hasse diagram has $e<s_1,s_2$, then $s_1,s_2<s_1s_2,s_2s_1$,
then $s_1s_2,s_2s_1<w_0$, so there are $19$ comparable pairs: $6$ with $d=0$,
$8$ with $d=1$, $4$ with $d=2$ and one ($e<w_0$) with $d=3$.

\begin{proposition}\label{prop:S3}
For $x\le w$ in $S_3$, $R_{x,w}$ depends only on $d=\ell(w)-\ell(x)$:
\[
R_{x,w}=1,\quad q-1,\quad (q-1)^2=q^2-2q+1,\quad q^3-2q^2+2q-1
\qquad (d=0,1,2,3).
\]
\end{proposition}

\begin{proof}
We apply~\eqref{eq:R}. For $d=1$, take the covering pair $e<s_1$ and $s=s_1$: $ws=e$
and $xs=s_1>e$, so $R_{e,s_1}=(q-1)R_{e,e}+qR_{s_1,e}=q-1$, because
$s_1\not\le e$. For $d=2$, take $e<s_1s_2$ and $s=s_2$: $ws=s_1$ and $xs=s_2>e$,
so $R_{e,s_1s_2}=(q-1)R_{e,s_1}+qR_{s_2,s_1}=(q-1)^2$, because $s_2\not\le s_1$.
For $d=3$, take $w_0=s_1s_2s_1$ and $s=s_1$: $ws=s_1s_2$ and $xs=s_1>e$, so
$R_{e,w_0}=(q-1)R_{e,s_1s_2}+qR_{s_1,s_1s_2}$. Here
$R_{s_1,s_1s_2}=(q-1)R_{s_1,s_1}+qR_{s_1s_2,s_1}=q-1$ (using $s=s_2$). Hence
$R_{e,w_0}=(q-1)^3+q(q-1)=q^3-2q^2+2q-1$. The remaining pairs are similar.
Lean checks all $36$ pairs and all choices of $s$ (Section~\ref{sec:lean}), and
\texttt{verify.py} recomputes everything from the Hecke algebra. It is also
classical that $R_{x,w}=q-1$ whenever $d=1$ and $R_{x,w}=(q-1)^2$ whenever $d=2$,
in every Coxeter group~\cite[Ch.~5]{bb}.
\end{proof}

\section{Refutation under every reading of ``max''}\label{sec:readings}

\begin{theorem}\label{thm:main}
Let $\rho$ be any of the following readings of ``$\max R_{x,w}$'':
(a) the largest coefficient; (b) the largest absolute value of a coefficient;
(c) the sum of the absolute values of the coefficients; (d) the leading
coefficient; (e) the degree; (f) the number of nonzero terms; (g) the value
$R_{x,w}(q_0)$ at a fixed point $q_0$ (any $q_0\in\mathbb C$); (h) the maximum
of the function $q\mapsto R_{x,w}(q)$ over $\mathbb Z$ or $\mathbb R$. Then the identity
$\rho(R_{x,w})=f(\ell(w)-\ell(x))$ fails for some pair $x<w$ of $S_3$.
\end{theorem}

\begin{proof}
By Proposition~\ref{prop:S3}, with pairs $e<s_1$ ($d=1$), $e<s_1s_2$ ($d=2$) and
$e<w_0$ ($d=3$):

\begin{center}
\begin{tabular}{lcccl}
\toprule
reading & $d=1$ & $d=2$ & $d=3$ & failure\\
\midrule
conjectured $f(d)$ & $1$ & $2$ & $6$ & \\
(a) largest coefficient & $1$ & $1$ & $2$ & $d=2$: $1\ne2$\\
(b) largest $|$coefficient$|$ & $1$ & $2$ & $2$ & $d=3$: $2\ne6$\\
(c) sum of $|$coefficients$|$ & $2$ & $4$ & $6$ & $d=1$: $2\ne1$\\
(d) leading coefficient & $1$ & $1$ & $1$ & $d=2$: $1\ne2$\\
(e) degree & $1$ & $2$ & $3$ & $d=3$: $3\ne6$\\
(f) number of nonzero terms & $2$ & $3$ & $4$ & $d=1$: $2\ne1$\\
(g) value at $q_0=2$ & $1$ & $1$ & $3$ & $d=2$: $1\ne2$\\
\bottomrule
\end{tabular}
\end{center}

For (g) with an arbitrary $q_0\in\mathbb C$: the case $d=1$ requires $q_0-1=1$, so
$q_0=2$, and then $d=2$ gives $(q_0-1)^2=1\ne2$. For (h): $R_{e,s_1}(q)=q-1$ is
unbounded above on $\mathbb Z$ (for instance $R_{e,s_1}(3)=2>1=f(1)$), so it has no
maximum, let alone the maximum $f(1)=1$. If ``max'' means the maximum over a
compact interval such as $[0,1]$, then $\max_{[0,1]}(q-1)=0\ne1$.
\end{proof}

\begin{remark}[The pair $x=w$]
For $x=w$ we have $d=0$, $R_{w,w}=1$ and $f(0)=2^0(1-2^0)=0$. So the non-strict
statement ``for all $x\le w$'' fails for every reading with $\rho(1)\ne0$, which
covers (a)--(d), (f), (g) and (h). Only the degree reading (e) survives $d=0$, and
it fails at $d=3$.
\end{remark}

\begin{remark}[Maximum over all pairs with the same $d$]
One might read ``$\max R_{x,w}$'' as the maximum of $\rho(R_{x,w})$ over all pairs
with the given length difference $d$. In $S_3$ all pairs with the same $d$ have
the same R-polynomial (Proposition~\ref{prop:S3}), so this aggregated reading
coincides with the per-pair one and fails as above. Over all Coxeter groups,
$R_{x,w}=q-1$ for $d=1$ and $(q-1)^2$ for $d=2$ always, so the failures at
$d\le2$ in the table persist. The degree is always $d$, so (e) fails at $d=3$.
For (b) at $d=3$, Dyer's formula~\cite{dyer} gives
$\widetilde R_{x,w}(t)=t^3+at$ with $a\in\{0,1\}$, hence
$R_{x,w}\in\{(q-1)^3,\ (q-1)^3+q(q-1)\}$, whose largest absolute coefficients are
$3$ and $2$, never $6$. These general-group facts are quoted from the literature
and are not formalized. \texttt{verify.py} also surveys $S_4$, where the
aggregated maxima are far from $f(d)$ for every reading.
\end{remark}

\begin{remark}[Other normalizations]
Some authors use $\varepsilon_x\varepsilon_w R_{x,w}=(-1)^dR_{x,w}$ or the
polynomial $\widetilde R_{x,w}$ with
$R_{x,w}(q)=q^{d/2}\widetilde R_{x,w}(q^{1/2}-q^{-1/2})$. For $S_3$ these are
$1-q,\ (q-1)^2,\ 1-2q+2q^2-q^3$ and $t,\ t^2,\ t^3+t$ ($d=1,2,3$). Every reading
(a)--(g) still fails, at $d=2$ or $d=3$ (and (g) for every $q_0$). This is checked
by \texttt{verify.py}. Readings (b), (c), (e), (f) are insensitive to the sign
change and to reversing the coefficient list.
\end{remark}

\begin{remark}[Scope]
The conjecture is false as a statement about KL theory, and in particular about
the symmetric groups. It is not claimed false in every individual Coxeter group:
in $W=\mathbb Z/2$ the only strict pair has $R=q-1$, which matches $f(1)=1$ under
several readings.
\end{remark}

\begin{proposition}[The literal recursion]\label{prop:literal}
Let $m$ be an involution of a set $W$ (e.g.\ $x\mapsto xs$ for a fixed simple
reflection $s$). Let $R_{x,w}$ ($x,w\in W$) be polynomials, or even formal power
series, over $\mathbb Z$ with $R_{x,w}=qR_{m(x),w}+R_{x,m(w)}$ for all $x,w$.
Then $R\equiv0$.
\end{proposition}

\begin{proof}
Write $R_{x,w}=\sum_k r_k(x,w)q^k$. Apply the recursion to $(m(x),w)$ and
$(m(x),m(w))$ and add. The terms $R_{m(x),w}+R_{m(x),m(w)}$ cancel, leaving
$q\,(R_{x,w}+R_{x,m(w)})=0$. Hence $r_k(x,w)=-r_k(x,m(w))$ for all $k$. Now
induct on $k$. The coefficient of $q^k$ in the recursion is
$r_k(x,w)=r_{k-1}(m(x),w)+r_k(x,m(w))$, where $r_{-1}=0$ and $r_{k-1}=0$ by
induction. So $r_k(x,w)=r_k(x,m(w))=-r_k(x,w)$, and $r_k(x,w)=0$.
\end{proof}

\section{Lean formalization}\label{sec:lean}

The Lean project \texttt{lean4/} (Lean 4.19.0, core library only, namespace
\texttt{KLR}) contains:
\begin{itemize}
\item \texttt{S3}: the six permutations as lists, with \texttt{S3\_spec}
(distinct, each a permutation of $[1,2,3]$). \texttt{inv} is the inversion count
(length) and \texttt{swapPos}/\texttt{rmul} are right multiplication by
transpositions and simple reflections.
\item \texttt{BruhatLe}: Bruhat order, defined inductively as the
reflexive--transitive closure of $x\to xt$ with $\ell(xt)>\ell(x)$.
\texttt{bruhatB} is a bounded chain search, and \texttt{bruhatB\_iff} proves
$\texttt{bruhatB}\ x\ w=\texttt{true}\iff\texttt{BruhatLe}\ x\ w$ for $x\in S_3$.
\texttt{bruhat\_count} counts $19$ comparable pairs, and \texttt{bruhat\_length}
shows that $x\le w$ implies $\ell(x)\le\ell(w)$.
\item Polynomials in $\mathbb Z[q]$ are normalized coefficient lists.
\texttt{IsRFamily R} is exactly~\eqref{eq:R}: vanishing off the Bruhat order,
$R_{w,w}=1$, and the recursion for \emph{every} simple $s$ with $ws<w$.
\texttt{R0\_isRFamily} shows that the computed family satisfies it, so the
recursion is consistent for every choice of descent. \texttt{RFamily\_unique}
shows that any two families satisfying it agree (by induction on $\ell(w)$).
So all statements are about \emph{the} R-polynomials.
\item \texttt{R\_classification} proves Proposition~\ref{prop:S3}, and
\texttt{R\_values} gives $R_{e,e},R_{e,s_1},R_{e,s_1s_2},R_{e,w_0}$.
\item \texttt{Formula v d} is the conjectured identity multiplied by
$2^{\lceil d/2\rceil}$: $v\cdot2^{c}=2^d(2^{c}-1)$ with
$c=\lceil d/2\rceil=\lfloor (d+1)/2\rfloor$. \texttt{formula\_values} shows
$f(0..4)=0,1,2,6,12$. \texttt{Claim}~$\rho$~\texttt{R} is ``for all $x\le w$ in
$S_3$, \texttt{Formula}$(\rho(R_{x,w}),d)$''. \texttt{ClaimStrict} restricts to
$x\ne w$, \texttt{ClaimAgg} is the aggregated reading, and
\texttt{ClaimMaxValue} is reading (h) over $\mathbb Z$.
\item \texttt{conjecture\_00000001102\_false}: for every \texttt{R} with
\texttt{IsRFamily R}, $\neg$\texttt{Claim}, $\neg$\texttt{ClaimStrict} and
$\neg$\texttt{ClaimAgg} hold for readings (a)--(f) and for (g) at every integer
$q_0$, and $\neg$\texttt{ClaimMaxValue} holds.
\texttt{conjecture\_00000001102\_false\_R0} instantiates this with the concrete
family. \texttt{not\_claim\_diag} gives the $x=w$ counterexample for every
$\rho$ with $\rho(1)\ne0$.
\item \texttt{literal\_recursion\_forces\_zero} is
Proposition~\ref{prop:literal} for coefficient sequences $\mathbb N\to\mathbb Z$
and any involution. \texttt{literal\_recursion\_inconsistent} applies it to
$S_3$ with $s=s_1$.
\end{itemize}
Not in Lean: evaluation at non-integer $q_0$ (the same two-line argument), the
general-Coxeter-group facts of the remarks, and the alternative normalizations.
\texttt{verify.py} checks the latter. There is no \texttt{sorry}, no
\texttt{native\_decide} and no added axioms. \texttt{\#print axioms} shows at most
\texttt{propext}, \texttt{Classical.choice} and \texttt{Quot.sound}.

\section{Reproduction}
\begin{verbatim}
cd lean4 && lake build       # Lean 4.19.0, core only
python3 verify.py            # standard library only
pdflatex report.tex
\end{verbatim}
\texttt{verify.py} uses a different method from the Lean code. It computes
$T_{w^{-1}}^{-1}$ in the Hecke algebra of $S_3$ (and $S_4$) with Laurent
polynomials and reads off $R_{x,w}$. It decides Bruhat order by the tableau
criterion and checks that $R_{x,w}\ne0$ exactly when $x\le w$. It then evaluates
every reading for $R$, $(-1)^dR$ and $\widetilde R$ against $f(d)$ computed with
exact fractions, and searches $q_0\in[-50,50]$ for evaluation points.

\begin{thebibliography}{9}
\bibitem{kl} D.~Kazhdan and G.~Lusztig, Representations of Coxeter groups and
Hecke algebras, \emph{Invent. Math.} 53 (1979), 165--184.
\bibitem{humphreys} J.~E.~Humphreys, \emph{Reflection Groups and Coxeter Groups},
Cambridge University Press, 1990.
\bibitem{bb} A.~Bj\"orner and F.~Brenti, \emph{Combinatorics of Coxeter Groups},
GTM 231, Springer, 2005.
\bibitem{dyer} M.~J.~Dyer, Hecke algebras and shellings of Bruhat intervals,
\emph{Compositio Math.} 89 (1993), 91--115.
\end{thebibliography}

\end{document}
